% 核心结论：有损线圈的并联谐振
% 图原型：纯几何／对比。物理网络按本题节点、端口和符号构建。
\documentclass[tikz,border=7pt]{standalone}
\usepackage{fontspec,xeCJK,amsmath,amssymb}
\IfFontExistsTF{Times New Roman}{\setmainfont{Times New Roman}}{\setmainfont{TeX Gyre Termes}}
\IfFontExistsTF{Heiti SC}{\setCJKmainfont[AutoFakeBold=1.5]{Heiti SC}}{\setCJKmainfont[AutoFakeBold=1.5]{FandolHei-Regular}}
\usetikzlibrary{arrows.meta,calc,angles,quotes}
\definecolor{ink}{HTML}{243745}
\definecolor{blue}{HTML}{275DAB}
\definecolor{green}{HTML}{177C63}
\definecolor{red}{HTML}{BD4935}
\definecolor{purple}{HTML}{715393}
\begin{document}
\begin{tikzpicture}[font=\fontsize{9.5}{12}\selectfont,text=ink,
 v/.style={line width=1.1pt,-{Stealth[length=2.4mm,width=1.5mm]}},
 axis/.style={ink!65,line width=.45pt,-{Stealth[length=1.7mm,width=1.1mm]}},
 aux/.style={ink!40,densely dashed,line width=.5pt},
 lab/.style={fill=white,inner sep=2pt},
 card/.style={draw=blue!25,fill=blue!3,rounded corners=3pt,text width=6cm,align=left,inner sep=10pt}]
\path[use as bounding box] (-8,-3.35) rectangle (8,3.65);
\node[font=\fontsize{10.5}{13}\selectfont\bfseries,text width=15cm,align=center] at (0.00000,3.22000) {M6-028　有损线圈的并联谐振};
\begin{scope}[shift={(-4,0)}]
\draw[ink,thick] (-2.50000,1.50000) -- (-2.50000,0.42000);
\draw[ink,thick] (-2.50000,-0.42000) -- (-2.50000,-1.50000);
\draw[ink,thick] (-2.5,0) circle (.42);
\node[] at (-2.50000,0.00000) {$\sim$};
\node[] at (-2.88000,1.60000) {$+$};
\node[] at (-2.88000,-1.50000) {$-$};
\node[] at (-3.12000,0.00000) {$u$};
\draw[ink,thick] (-2.50000,1.50000) -- (2.25000,1.50000);
\draw[ink,thick] (-2.50000,-1.50000) -- (2.25000,-1.50000);
\node[] at (-0.90000,1.80000) {$A$};
\node[] at (-1.70000,-1.80000) {$B$};
\draw[ink,thick] (-0.45000,1.50000) -- (-0.45000,0.98000);
\begin{scope}[shift={(-0.45000,0.50000)},rotate=-90]
\draw[ink,thick] (-.48,-.16) rectangle (.48,.16);
\end{scope}
\draw[ink,thick] (-0.45000,0.02000) -- (-0.45000,-0.02000);
\begin{scope}[shift={(-0.45000,-0.50000)},rotate=-90]
\draw[ink,thick] plot[domain=-.48:.48,samples=80] (\x,{.22*abs(sin(750*(\x+.48)))});
\end{scope}
\draw[ink,thick] (-0.45000,-0.98000) -- (-0.45000,-1.50000);
\node[] at (-1.00000,0.50000) {$R$};
\node[] at (-1.00000,-0.50000) {$L$};
\draw[v,blue] (0.10000,1.05000) -- (0.10000,0.35000);
\node[lab,text=blue,anchor=west] at (0.25000,0.70000) {$i_L$};
\draw[ink,thick] (1.75000,1.50000) -- (1.75000,0.48000);
\begin{scope}[shift={(1.75000,0.00000)},rotate=-90]
\draw[ink,thick] (-.10,-.35) -- (-.10,.35);
\draw[ink,thick] (.10,-.35) -- (.10,.35);
\draw[ink,thick] (-.48,0) -- (-.10,0);
\draw[ink,thick] (.10,0) -- (.48,0);
\end{scope}
\draw[ink,thick] (1.75000,-0.48000) -- (1.75000,-1.50000);
\node[] at (1.20000,0.00000) {$C$};
\draw[v,green] (2.15000,1.05000) -- (2.15000,0.35000);
\node[lab,text=green,anchor=west] at (2.20000,0.70000) {$i_C$};
\draw[v,red] (-2.20000,1.50000) -- (-1.45000,1.50000);
\node[lab,text=red,anchor=center] at (-1.82500,1.82000) {$i$};
\node[] at (0.00000,-2.40000) {$u=u_{AB}$；各支路电流从 $A$ 流向 $B$。};
\end{scope}
\begin{scope}[shift={(4,0)}]
\node[card] at (0.00000,1.15000) {$Y=1/(R+j\omega L)+j\omega C$};
\end{scope}
\begin{scope}[shift={(4,0)}]
\node[card] at (0.00000,-1.15000) {$\omega_0^2=1/(LC)-(R/L)^2$};
\end{scope}
\node[text width=15cm,align=center] at (0.00000,-3.00000) {图中的R与L在同一支路串联；电容补偿该支路电流的无功分量。};

\end{tikzpicture}\end{document}